% GENERATED FILE. Edit canonical lessons, then run npm run generate:books.
% canonical-source-hash: fnv1a64:1913b8d77237eac1
% canonical-lessons: SA-C445-ulka, SA-C445-dhruvah, SA-C445-grahanam, SA-C445-diva, SA-C445-yapayati

\chapter{Stars by Night and Day}
\label{ch:sa-stars-by-night-and-day}
\hlchaptermodality{445}

\begin{chapteropening}
\textbf{By the end of this chapter:} \emph{I can talk about shooting stars, the pole star and eclipses.}

Complete the last lesson of chapter 445: I can talk about shooting stars, the pole star and eclipses.
\end{chapteropening}

\section[\texorpdfstring{\sk{उल्का} (ulkā)}{उल्का (ulkā)}]{\sk{उल्का} (ulkā) — a shooting star; a meteor}

\label{lesson:SA-C445-ulka}



\begin{quote}
Before the new one: say the Sanskrit for the dark fortnight, then the Sanskrit for gleams.
\end{quote}

\subsection*{You'll want to know: \sk{उल्का}}
\textbf{\sk{उल्का}} — \emph{ulkā} — \textquotedblleft{}a shooting star; a meteor\textquotedblright{}.

\begin{quote}
\sk{रात्रौ} \sk{आकाशे} \sk{उल्का} \sk{पतति।}
\end{quote}

\begin{quote}
\emph{rātrau ākāśe ulkā patati.}
\end{quote}

\begin{quote}
At night a shooting star falls across the sky.
\end{quote}

A feminine noun.

\begin{cousinweb}
\textbf{\sk{उल्का}} (\emph{ulkā}) comes from the root √uṣ, \textquotedblleft{}to burn\textquotedblright{}, to which MW refers it; MW gives \emph{ulkā} as a fiery phenomenon in the sky, a meteor, a firebrand (MW).
\end{cousinweb}

\begin{grammarlens}[title={What you've built}]
One more word for getting around.
\end{grammarlens}

\subsection*{Your turn}
\begin{itemize}
\raggedright
  \item \emph{Say it:} \emph{ulkā}
  \item \emph{Say it:} \emph{ulkā}, once more
  \item \emph{Say it:} \emph{rātrau ākāśe ulkā patati.} — At night a shooting star falls across the sky.
  \item \emph{Recall:} say the Sanskrit for the dark fortnight, then the Sanskrit for gleams, then say \emph{ulkā} again
\end{itemize}

\begin{tcolorbox}[breakable,colback=teal!4,colframe=teal!35!black,title={Before you move on}]
What does \sk{उल्का} mean? (\textquotedblleft{}A shooting star; a meteor\textquotedblright{}.) Say it once more.
\end{tcolorbox}
\section[\texorpdfstring{\sk{ध्रुवः} (dhruvaḥ)}{ध्रुवः (dhruvaḥ)}]{\sk{ध्रुवः} (dhruvaḥ) — the pole star}

\label{lesson:SA-C445-dhruvah}



\begin{quote}
Before the new one: say the Sanskrit for gleams, then the Sanskrit for a shooting star.
\end{quote}

\subsection*{You'll want to know: \sk{ध्रुवः}}
\textbf{\sk{ध्रुवः}} — \emph{dhruvaḥ} — \textquotedblleft{}the pole star\textquotedblright{}.

\begin{quote}
\sk{ध्रुवः} \sk{उत्तरायां} \sk{दिशायां} \sk{तिष्ठति।}
\end{quote}

\begin{quote}
\emph{dhruvaḥ uttarāyāṃ diśāyāṃ tiṣṭhati.}
\end{quote}

\begin{quote}
The pole star stands in the north.
\end{quote}

A masculine noun.

\begin{cousinweb}
\textbf{\sk{ध्रुवः}} (\emph{dhruvaḥ}) comes from the root √dhṛ, \textquotedblleft{}to hold, to bear\textquotedblright{}, as MW gives it: MW gives \emph{dhruva} as fixed, firm, and as the pole star, the star that does not move (MW).
\end{cousinweb}

\begin{grammarlens}[title={What you've built}]
One more word for getting around.
\end{grammarlens}

\subsection*{Your turn}
\begin{itemize}
\raggedright
  \item \emph{Say it:} \emph{dhruvaḥ}
  \item \emph{Say it:} \emph{dhruvaḥ}, once more
  \item \emph{Say it:} \emph{dhruvaḥ uttarāyāṃ diśāyāṃ tiṣṭhati.} — The pole star stands in the north.
  \item \emph{Recall:} say the Sanskrit for gleams, then the Sanskrit for a shooting star, then say \emph{dhruvaḥ} again
\end{itemize}

\begin{tcolorbox}[breakable,colback=teal!4,colframe=teal!35!black,title={Before you move on}]
What does \sk{ध्रुवः} mean? (\textquotedblleft{}The pole star\textquotedblright{}.) Say it once more.
\end{tcolorbox}
\section[\texorpdfstring{\sk{ग्रहणम्} (grahaṇam)}{ग्रहणम् (grahaṇam)}]{\sk{ग्रहणम्} (grahaṇam) — an eclipse}

\label{lesson:SA-C445-grahanam}



\begin{quote}
Before the new one: say the Sanskrit for a shooting star, then the Sanskrit for the pole star.
\end{quote}

\subsection*{You'll want to know: \sk{ग्रहणम्}}
\textbf{\sk{ग्रहणम्}} — \emph{grahaṇam} — \textquotedblleft{}an eclipse\textquotedblright{}.

\begin{quote}
\sk{अद्य} \sk{सूर्यस्य} \sk{ग्रहणं} \sk{भविष्यति।}
\end{quote}

\begin{quote}
\emph{adya sūryasya grahaṇaṃ bhaviṣyati.}
\end{quote}

\begin{quote}
Today there will be an eclipse of the sun.
\end{quote}

A neuter noun.

\begin{cousinweb}
\textbf{\sk{ग्रहणम्}} (\emph{grahaṇam}) is built on the root √grah, \textquotedblleft{}to seize\textquotedblright{}, which also gives \emph{gṛhṇāti}, \textquotedblleft{}takes\textquotedblright{}: a seizing; MW gives \emph{grahaṇa} as seizing and as an eclipse, the sun or moon being seized (MW).
\end{cousinweb}

\begin{grammarlens}[title={What you've built}]
One more word for getting around.
\end{grammarlens}

\subsection*{Your turn}
\begin{itemize}
\raggedright
  \item \emph{Say it:} \emph{grahaṇam}
  \item \emph{Say it:} \emph{grahaṇam}, once more
  \item \emph{Say it:} \emph{adya sūryasya grahaṇaṃ bhaviṣyati.} — Today there will be an eclipse of the sun.
  \item \emph{Recall:} say the Sanskrit for a shooting star, then the Sanskrit for the pole star, then say \emph{grahaṇam} again
\end{itemize}

\begin{tcolorbox}[breakable,colback=teal!4,colframe=teal!35!black,title={Before you move on}]
What does \sk{ग्रहणम्} mean? (\textquotedblleft{}An eclipse\textquotedblright{}.) Say it once more.
\end{tcolorbox}
\section[\texorpdfstring{\sk{दिवा} (divā)}{दिवा (divā)}]{\sk{दिवा} (divā) — by day}

\label{lesson:SA-C445-diva}



\begin{quote}
Before the new one: say the Sanskrit for the pole star, then the Sanskrit for an eclipse.
\end{quote}

\subsection*{You'll want to know: \sk{दिवा}}
\textbf{\sk{दिवा}} — \emph{divā} — \textquotedblleft{}by day\textquotedblright{}.

\begin{quote}
\sk{दिवा} \sk{आकाशे} \sk{सूर्यः} \sk{अस्ति}, \sk{नक्तं} \sk{चन्द्रः।}
\end{quote}

\begin{quote}
\emph{divā ākāśe sūryaḥ asti, naktaṃ candraḥ.}
\end{quote}

\begin{quote}
By day the sun is in the sky, by night the moon.
\end{quote}

\begin{cousinweb}
\textbf{\sk{दिवा}} (\emph{divā}) comes from \emph{div}, \textquotedblleft{}the sky, day\textquotedblright{}, of which MW gives it as an old case form; the same word gives \emph{divya}, \textquotedblleft{}divine\textquotedblright{}; MW gives \emph{divā} as by day, often set against \emph{naktam}, \textquotedblleft{}by night\textquotedblright{} (MW).
\end{cousinweb}

\begin{grammarlens}[title={What you've built}]
One more word for getting around.
\end{grammarlens}

\subsection*{Your turn}
\begin{itemize}
\raggedright
  \item \emph{Say it:} \emph{divā}
  \item \emph{Say it:} \emph{divā}, once more
  \item \emph{Say it:} \emph{divā ākāśe sūryaḥ asti, naktaṃ candraḥ.} — By day the sun is in the sky, by night the moon.
  \item \emph{Recall:} say the Sanskrit for the pole star, then the Sanskrit for an eclipse, then say \emph{divā} again
\end{itemize}

\begin{tcolorbox}[breakable,colback=teal!4,colframe=teal!35!black,title={Before you move on}]
What does \sk{दिवा} mean? (\textquotedblleft{}By day\textquotedblright{}.) Say it once more.
\end{tcolorbox}
\section[\texorpdfstring{\sk{यापयति} (yāpayati)}{यापयति (yāpayati)}]{\sk{यापयति} (yāpayati) — passes the time}

\label{lesson:SA-C445-yapayati}



\begin{quote}
Before the new one: say the Sanskrit for an eclipse, then the Sanskrit for by day.
\end{quote}

\subsection*{You'll want to know: \sk{यापयति}}
\textbf{\sk{यापयति}} — \emph{yāpayati} — \textquotedblleft{}passes the time\textquotedblright{}.

\begin{quote}
\sk{वयं} \sk{प्रपायां} \sk{रात्रिं} \sk{यापयामः।}
\end{quote}

\begin{quote}
\emph{vayaṃ prapāyāṃ rātriṃ yāpayāmaḥ.}
\end{quote}

\begin{quote}
We pass the night at the water shed.
\end{quote}

Past (he …): \sk{सः} \sk{तत्र} \sk{रात्रिं} \sk{यापितवान्}, \emph{saḥ tatra rātriṃ yāpitavān}.

\begin{cousinweb}
\textbf{\sk{यापयति}} (\emph{yāpayati}) is built on the root √yā, \textquotedblleft{}to go\textquotedblright{}, which gives \emph{yāti}, \textquotedblleft{}travels\textquotedblright{}, in its causative, \textquotedblleft{}makes go\textquotedblright{}: to let time go by, to pass it (MW).
\end{cousinweb}

\begin{grammarlens}[title={What you've built}]
One more word for getting around.
\end{grammarlens}

\subsection*{Your turn}
\begin{itemize}
\raggedright
  \item \emph{Say it:} \emph{yāpayati}
  \item \emph{Say it:} \emph{yāpayati}, once more
  \item \emph{Say it:} \emph{vayaṃ prapāyāṃ rātriṃ yāpayāmaḥ.} — We pass the night at the water shed.
  \item \emph{Recall:} say the Sanskrit for an eclipse, then the Sanskrit for by day, then say \emph{yāpayati} again
\end{itemize}

\begin{tcolorbox}[breakable,colback=teal!4,colframe=teal!35!black,title={Before you move on}]
What does \sk{यापयति} mean? (\textquotedblleft{}Passes the time\textquotedblright{}.) Say it once more.
\end{tcolorbox}
